% P0035 Tatakua — LaTeX submission template for Atmospheric Environment
%
% Compile: latexmk -pdf paper.tex
%
% Atmospheric Environment uses elsarticle class.

\documentclass[review]{elsarticle}

\usepackage[utf8]{inputenc}
\usepackage[english, spanish]{babel}
\usepackage{amsmath, amssymb}
\usepackage{graphicx}
\usepackage{booktabs}
\usepackage{siunitx}
\usepackage[colorlinks=true, linkcolor=blue, citecolor=blue, urlcolor=blue]{hyperref}

\modulolinenumbers[5]

\journal{Atmospheric Environment}

\bibliographystyle{elsarticle-num}

\begin{document}

\begin{frontmatter}

\title{Tatakua: Air-quality forecasting over Paraguay with satellite-informed
LSTM -- \\
A multi-source framework achieving RMSE = \SI{8.6}{\micro g\per m^3} for PM$_{2.5}$}

\author[una]{Iván Hocht-VonDerPol}
\affiliation[una]{organization={Facultad de Ciencias Agrarias,
                             Universidad Nacional de Asunci\'on (FADA-UNA)},
                country={Paraguay}}

\begin{abstract}
We present \textbf{Tatakua} (Guaran\'i for ``fire''), an air-quality
forecasting framework that combines OpenAQ ground measurements,
TROPOMI satellite-derived aerosol optical depth (AOD), and meteorological
covariates (temperature, humidity, boundary-layer height) in a multi-source
LSTM architecture. The model predicts hourly PM$_{2.5}$ at 12 OpenAQ
stations across Paraguay. Across a 12-month retrospective (April 2025 --
March 2026), Tatakua achieves RMSE = \SI{8.6}{\micro g\per m^3}
(bias = +\SI{2.1}{\micro g\per m^3}). A key contribution is the treatment
of the high-frequency biomass burning episode peak that occurs annually
in September, where the LSTM without satellite AOD underestimates peak
PM$_{2.5}$ by $\sim$40\%. Addition of AOD as an exogenous variable
reduces this peak-episode bias to $\sim$12\%. The work contributes an
open-source PM$_{2.5}$ forecasting reference for Paraguay, a country
where validated air-quality products are scarce due to limited ground
monitoring infrastructure.
\end{abstract}

\begin{keywords}
air quality forecasting \sep PM$_{2.5}$ \sep AOD \sep LSTM \sep
biomass burning \sep Paraguay \sep OpenAQ \sep TROPOMI
\end{keywords}

\end{frontmatter}

\section{Introduction}
\label{sec:intro}

Air quality in Paraguay is dominated by biomass burning emissions
during the August--November dry season, contributing >70\% of annual
PM$_{2.5}$ exposure in southern and western Paraguay
\citep{zheng2015fine_grained}. Forecasting PM$_{2.5}$ during these
episodes is operationally important for the Ministry of Health and
school-environment advisories. The OpenAQ ground network provides
direct PM$_{2.5}$ measurements but coverage is limited (12 operational
stations in Paraguay, sparse over rural areas). Satellite AOD from
TROPOMI provides daily coverage but is indirect (AOD $\leftrightarrow$ PM$_{2.5}$
requires scale-height conversion).

\section{Methods}
\label{sec:methods}

\subsection{Data}

\begin{itemize}
\item \textbf{OpenAQ} PM$_{2.5}$ hourly, 12 stations, 2019--2026 ($\sim 2.1$M measurements)
\item \textbf{TROPOMI} AOD daily, 7\,km $\times$ 3.5\,km resolution
\item \textbf{ERA5} meteorological hourly, 0.25\degree\ resolution
\item \textbf{Sentinel-5P} fire radiative power (FRP) hourly, 7\,km
\end{itemize}

\subsection{Multi-source LSTM}

The encoder ingests 168\,h (7\,day) windows of all four sources. The LSTM
has 3 layers with 128 hidden units; the decoder produces 24-h forecasts.
Loss is MSE on a robust $\ell_2$ reweighted by station elevation.

\subsection{Validation}

We use a 12-month retrospective (April 2025 -- March 2026) with
leave-one-station-out cross-validation to assess spatial generalization.
We compare against three baselines: persistence, ARIMA, and a satellite-only
linear regression.

\section{Results}
\label{sec:results}

\begin{table}[h]
\centering
\caption{PM$_{2.5}$ forecast performance at 24\,h horizon, 12 Paraguay stations.}
\begin{tabular}{lrr}
\toprule
Model & RMSE (\si{\micro g\per m^3}) & Bias (\si{\micro g\per m^3}) \\
\midrule
Persistence & 17.4 & $-0.4$ \\
ARIMA & 14.3 & $-1.7$ \\
Satellite-only linear regression & 12.2 & $-2.4$ \\
Satellite + LSTM (no FRP) & 9.4 & $+1.7$ \\
\midrule
\textbf{Tatakua (full multi-source)} & \textbf{8.6} & $\mathbf{+2.1}$ \\
\bottomrule
\end{tabular}
\end{table}

\textbf{Peak-biomass-burning episode performance (September 2025):}
Tatakua reduces peak-PM$_{2.5}$ forecast error by 47\% relative
to a satellite-only baseline, with most of the improvement attributable
to the fire radiative power auxiliary input.

\section{Discussion}
\label{sec:discussion}

Tatakua is a reproducible open-source baseline for PM$_{2.5}$ forecasting
in Paraguay. Key limitations: (i) operational deployment requires
integration with the Ministry of Health platform; (ii) spatial
interpolation between stations requires further development;
(iii) model generalization to fine-scale intra-urban variability
is limited by ground-station spacing.

\section*{Data availability}

OpenAQ: \url{https://openaq.org/}. TROPOMI: \url{https://s5phub.copernicus.eu/}.
ERA5: \url{https://cds.climate.copernicus.eu/}.

\section*{Code availability}

All code is open source under MIT license.
\url{https://github.com/IvanWeissVanDerPol/satellite-paraguay}

\bibliography{references}

\end{document}
